\scanpage{123}
\mspage{38}
% archive sheet number: 122

% L01-02
Soit
\[
U=X-\operatorname{supp}D,
\]

% L03-04
et supposons que
\[
\deleted{\unsure{a}}t|U\in\Gamma(U,\unsure{\mathcal{O}|\mathscr{O}}_U^*),\quad\text{i.e. }\operatorname{supp}\operatorname{div}t\subset\operatorname{supp}D.
\]

% L05-08
\deleted{Posons, pour tout entier $m$ premier aux caractéristiques} Supposons pour simplifier \deleted{\unsure{qu'il existe} \unsure{$X$ est de car nulle}, \inserted{$X$, de car. nulle, \unsure{ou bien}} \unsure{qu'il existe au plus un nb premier} $p$ \unsure{qui ne soit inversible sur} $X$.} $V(t)$ soit de car. $p>0$ et que tt nb premier $\ell\ne p$ soit inversible sur $X$.

% L09-10
Pour tout \deleted{\unsure{nb premier}} entier $m>0$ premier aux car. résiduelles de $X$, \deleted{\unsure{posons}} soit

% L11-12
\[
\left\{\begin{aligned}
X'(m)&=\operatorname{Spec}\mathcal{O}_X[T]/(T^m-t)\\
U'(m)&=X'(m)\times_X U
\end{aligned}\right.,
\]

% L13-15
de sorte que les $U'(m)$ forment un système projectif de rev. \deleted{étales et} principaux de $U$, de groupes les $\mu_{m,U}$. On considère

% L16
\[
U'=\varprojlim U'(m),
\]

% L17-18
qui est principal sur $U$ de groupe \deleted{$\unsure{\mu_m}$} $G_U$, où
\[
G=\varprojlim\mu_m\deleted{\unsure{(X)}},
\]

% L19-20
et on \deleted{\unsure{suppose}} considère \unsure{le|la} \deleted{projection}\inserted{\unsure{composé|composée}}
\diagram{\begin{tikzcd}[ampersand replacement=\&]
U'\arrow[r]\arrow[rr,bend left=22,"f"]\&U\arrow[r]\&X.
\end{tikzcd}}

% L21-24
Soit $n\geq1$ premier aux car. résiduelles de $X$, et \deleted{$\Lambda=\unsure{\mathbb{Z}|Z}/n\mathbb{Z}$} $\Lambda=\mathbb{Z}/n\mathbb{Z}$, on s'intéresse aux faisceaux
\[
R^i f_*(\Lambda_{U'})|X_0,\qquad X_0=V(t)
\]
